\documentclass{article}

\usepackage[utf8]{inputenc}
\usepackage{amssymb,amsmath}% 
\usepackage{pifont}
\usepackage{graphicx}
\usepackage{verbatimbox}
\usepackage{hyperref} 
\usepackage{tcolorbox}

\usepackage{xspace}
\usepackage{booktabs}
\usepackage{multirow}
\usepackage{rotating}
\usepackage{color}
\usepackage{palatino}
\usepackage[title]{appendix}
\usepackage{sectsty}
\usepackage{makecell}

\newcommand{\eg}{\emph{e.g.}}
\newcommand{\etal}{\emph{et.al.}}
\newcommand{\ie}{\emph{i.e.}}
\newcommand{\etc}{\emph{etc}}
\newcommand{\wrt}{\emph{w.r.t.}}

%
\title{\Large\textbf{ParameterNet: Parameters Are All You Need for Large-scale Visual Pretraining of Mobile Networks}}

%
\author{
\begin{minipage}{\linewidth}
\begin{center}
\normalsize Kai Han$^{1,2,*}$ \hspace{0.23cm} Yunhe Wang$^{1,}$\footnote{Equal contribution.} \hspace{0.23cm} Jianyuan Guo$^{1}$ \hspace{0.23cm} Enhua Wu$^{2,3}$ \\[0.2cm] 
{\normalsize$^1$Huawei Noah's Ark Lab}\\
{\normalsize$^2$State Key Lab of Computer Science, ISCAS \& UCAS}\\
{\normalsize$^3$University of Macau}\\
{\small\texttt{\{kai.han,yunhe.wang,jianyuan.guo\}@huawei.com, weh@ios.ac.cn}}\\
\end{center}
\end{minipage}
}

\date{~}

\begin{document}

\maketitle

\begin{abstract}
	The large-scale visual pretraining has significantly improve the performance of large vision models. However, we observe the \emph{low FLOPs pitfall} that the existing low-FLOPs models cannot benefit from large-scale pretraining.
	In this paper, we propose a general design principle of adding more parameters while maintaining low FLOPs for large-scale visual pretraining, named as ParameterNet. Dynamic convolutions are used for instance to equip the networks with more parameters and only slightly increase the FLOPs. The proposed ParameterNet scheme enables low-FLOPs networks to benefit from large-scale visual pretraining. Experiments on the large-scale ImageNet-22K have shown the superiority of our ParameterNet scheme. For example, ParameterNet-600M can achieve higher accuracy than the widely-used Swin Transformer (81.6\% \emph{vs.} 80.9\%) and has much lower FLOPs (0.6G \emph{vs.} 4.5G). The code will be released as soon\footnote{MindSpore: \url{https://gitee.com/mindspore/models}}\footnote{PyTorch: \url{https://github.com/huawei-noah/Efficient-AI-Backbones}}.
\end{abstract}

%%%%%%%%% BODY TEXT
\section{Introduction}
\label{sec:intro}
Thanks to the development of computation hardware and data engineering, the large-scale visual pretraining as a fundamental component in computer vision has made impressive progress. The pretrained vision models serve as the representation learners and transfer to downstream tasks such as image recognition~\cite{vit,efficientnetv2} and object detection~\cite{swin}.

The mainstream pretrained vision models usually requires a large amount of resources including data, parameters and FLOPs. These three key factors heavily influence the performance and basically follow the scaling law~\cite{zhai2022scaling}. The large pretraining data can provide diverse samples for representation learning. The sizes of these datasets range from millions~\cite{imagenet,kuznetsova2020open} to billions~\cite{sun2017revisiting,thomee2015new}, for example, the widely-used ImageNet-22K dataset~\cite{imagenet} consists of 14M images and 21,841 categories. To better fitting on the large dataset, the model sizes (including both parameters and FLOPs) are getting larger and larger in recent years, \eg, ViT-G/14 model has 1.8B parameters and 965B FLOPs~\cite{zhai2022scaling}.

The visual applications on mobile devices usually requires fast inference, so it is difficult to deploy the existing pretrained vision models due to the high computational cost. To address this issue, we empirically study the effect of FLOPs in large-scale visual pretraining. ImageNet-22K is adopted as the large-scale pretraining data and ImageNet-1K is a relatively small dataset for comparison. 
The pretrained transformer and CNN models are then finetuned on ImageNet-1K to evaluate the performance. As shown in Figure~\ref{fig:swin} and \ref{fig:enetv2}, when model FLOPs gradually increase, the model accuracy increases consistently. For the high-FLOPs models, 22K pretrained models outperform 1K ones. However, the low-FLOPs models cannot benifit from large-scale pretraining, and we called this observation as \emph{low FLOPs pitfall}. 

In this paper, we construct low-FLOPs ParameterNet by adding more parameters while maintaining low FLOPs for large-scale visual pretraining. It is a general design principle and there are various approaches with more parameters and low FLOPs. For instance, here we mainly consider the efficient dynamic convolution which manyfold increase the number of parameters while almost does not bring in extra FLOPs. The ParameterNet scheme can enable the previous networks to benefit from the large-scale visual pretraining and overcome the \emph{low FLOPs pitfall}. In the experiments, the ImageNet-22K pretrained ParameterNets can improve the performance by about +2\% over the regular ImageNet-1K training. For example, ParameterNet-600M achieves 81.6\% top-1 accuracy on ImageNet-1K val set whose \#FLOPs is 7$\times$ lower than that of Swin-T.

\begin{figure}[htp]
	\centering
	\includegraphics[width=1.0\linewidth]{figures/ghostnet-pnet.png}
	\caption{Results on ImageNet-1K validation set. The original GhostNet falls into the \emph{low FLOPs pitfall}. The proposed ParameterNet overcomes the \emph{low FLOPs pitfall}.}
	\label{fig:main}
\end{figure}

\section{Related Work}
\subsection{Visual Backbone Networks}
The deep neural networks in computer vision can be divided into CNNs, vision transformers and others. CNN used to be the mainstream network architecture for visual tasks~\cite{lecun1998gradient,alexnet,resnet}. The first trainable CNN \ie, LeNet~\cite{lecun1998gradient} is applied on optimal character recognition (a typical visual task). From 2012, CNNs began to be deeper and larger for more complex visual tasks, such as image classification~\cite{alexnet}, object detection~\cite{fasterRCNN} and semantic segmentation~\cite{fcn}. ResNet~\cite{resnet} introduces the shortcut connection to train the deeper networks and is widely used in vision and other communities. MobileNet~\cite{mobilenet} is designed for mobile devices and EfficientNet~\cite{efficientnet} scales the network from small to large.

Vision transformer is introduced into visual tasks from 2020~\cite{vit-survey,vit,detr}. ViT~\cite{vit} is the first transformer backbone by dividing the image into patches and processing them using the standard transformer architcture. Then a number of variants and improvements are proposed incluidng the pyramid architectures~\cite{pvt,swin}, the local attentions~\cite{tnt,swin} and hybrid networks~\cite{guo2022cmt,metaformer}.

Beyond CNNs and transformers, other types of neural networks are also explored for visual tasks. MLP-like architectures~\cite{mixer,resmlp} with only fully-connected layers as main operators can potentially simplify the software and hardware design for AI. The improved versions of MLP~\cite{cyclemlp,asmlp,hire,wavemlp} can enhance locality and translation equivalence. GNN has also been expored in vision and achieves competitive performance to transformers~\cite{vig}.

\subsection{Visual Pretraining}
Large-scale pretraining has achieved great success on natural language processing such as GPT series~\cite{gpt3,gpt4}. In the field of computer vision, large-scale pretraining is also beneficial and helps for downstream tasks~\cite{bit,swin,sam}.
The large datasets are the foundations of pretraining. To distinguish from the regular ImageNet-1K training, we consider the dataset with an order of magnitude more than ImageNet-1K (\ie, more than 10M samples) as the large-scale dataset. The commonly-used large visual datasets include ImageNet-22K~\cite{imagenet}, JFT-300M~\cite{sun2017revisiting}, YFFCC100M~\cite{yfcc100m} and IG-1B-Targeted~\cite{yalniz2019billion}.
The supervised pretraining is popular as it can learn semantic and meaningful representations for downstream visual tasks like segmentation and detection. BiT~\cite{bit} pretrained on JFT-300M achieves state-of-the-art on fine-grained recognition datasets. The ImageNet-22K pretrained Swin Transformer~\cite{swin} obtains high performance on segmentation and detection tasks.
The unsupervised pretraining appeals many researchers as it does not requires labels and may leverage the massive unlabeled data. The most popular approaches of visual unsupervised pretraining are contrastive learning~\cite{he2020momentum,simclr} and masked image modeling (MIM)~\cite{bao2021beit,he2022masked,xie2022simmim}.
In this paper, we utilize the vanilla supervised pretraining for simplicity and generality.

\section{Low FLOPs Pitfall}
The computational cost (\ie, FLOPs) is an important term in the scaling of visual models. We first investigate the effect of FLOPs and observe inspiring phenomenon. Both transformer and CNN architectures are studied on ImageNet-22K and ImageNet-1K pretraining.

\paragraph{Transformer.}
Swin Transformer~\cite{swin} is a representative vision transformer architecture with window attention and shifted window. We reproduce the models using the official code\footnote{https://github.com/microsoft/Swin-Transformer} and pretrain Swin Transformers with different scales on both ImageNet-22K and ImageNet-1K. The ImageNet-1k finetuning results are reported in Figure~\ref{fig:swin} for comparison. From the results, we can see that the accuracy increases as the FLOPs increase gradually with both ImageNet-1K and ImageNet-22K pretraining. For models with high FLOPs ($>$10G), pretraining on ImageNet-22K outperforms that on ImageNet-1K. However, pretraining on more data does not improve the performance for models with lower FLOPs ($<$4G).

\begin{figure}[htp]
	\centering
	\includegraphics[width=0.5\linewidth]{figures/swin-acc-flops.png}
	\renewcommand\arraystretch{1.05}
	\resizebox{.45\columnwidth}{!}{
		\begin{tabular}[b]{l|cccc}
			\hline
			Model & Pretrain data & \#Params & \#FLOPs & Top1 \\
			\hline
			Swin-300M & IN-1K & 1.9M & 0.3G & 69.2 \\
			Swin-T & IN-1K & 28M & 4.5G & 81.2 \\
			Swin-S & IN-1K & 50M & 8.7G & 83.2 \\
			Swin-B & IN-1K & 88M & 15.4G & 83.5 \\
			Swin-L & IN-1K & 197M & 34.5G & 83.5 \\
			\hline
			Swin-300M & IN-22K & 1.9M & 0.3G & 67.0 \\
			Swin-T & IN-22K & 28M & 4.5G & 80.9 \\
			Swin-S & IN-22K & 50M & 8.7G & 83.2 \\
			Swin-B & IN-22K & 88M & 15.4G & 85.2 \\
			Swin-L & IN-22K & 197M & 34.5G & 86.3 \\
			\hline
			\multicolumn{5}{c}{\vspace{2.5em}}\\
	\end{tabular}}
	\caption{Low FLOPs pitfall. Swin Transformer results on ImageNet-1K validation set. The red and blue lines denote ImageNet-21K and ImageNet-1K pretraining, respectively.}
	\label{fig:swin}
\end{figure}

\paragraph{CNN.}
For CNN, we select the widely-used EfficientNetV2~\cite{efficientnetv2} which is a family of convolutional networks scaling from small to large. We use the official code\footnote{https://github.com/google/automl/tree/master/efficientnetv2} and pretrain the models on both ImageNet-22K and ImageNet-1K. From the ImageNet-1k finetuning results in Figure~\ref{fig:enetv2}, we can observe the similar trend as that in Swin Transformer, especially, EfficientNetV2 models with less than 2G FLOPs pretraining on ImageNet-22K cannot perform better than those pretraining on ImageNet-1K.

From the observations of both transformer and CNN networks, we have a empirical conclusion that low-FLOPs models cannot benefit from large-scale pretraining, which is named as \emph{low FLOPs pitfall}.

\begin{figure}[htp]
	\centering
	% \fbox{\rule{0pt}{2in} \rule{0.9\linewidth}{0pt}}
	\includegraphics[width=0.5\linewidth]{figures/enetv2-acc-flops.png}
	\renewcommand\arraystretch{1.05}
	\resizebox{.45\columnwidth}{!}{
		\begin{tabular}[b]{l|cccc}
			\hline
			Model & Pretrain data & \#Params & \#FLOPs & Top1 \\
			\hline
			EfficientNetV2-B0 & IN-1K & 7.1M & 0.72G & 78.7 \\
			EfficientNetV2-B1 & IN-1K & 8.1M & 1.2G & 79.8 \\
			EfficientNetV2-B2 & IN-1K & 10.1M & 1.7G & 80.5 \\
			EfficientNetV2-B3 & IN-1K & 14.4M & 3.0G & 82.1 \\
			EfficientNetV2-S & IN-1K & 21.5M & 8.4G & 83.9 \\
			EfficientNetV2-M & IN-1K & 54.1M & 24.7G & 85.2 \\
			%EfficientNetV2-L & IN-1K & 119.5M & 56.3G & 85.7 \\
			\hline
			EfficientNetV2-B0 & IN-22K & 7.1M & 0.72G & 77.6 \\
			EfficientNetV2-B1 & IN-22K & 8.1M & 1.2G & 79.0 \\
			EfficientNetV2-B2 & IN-22K & 10.1M & 1.7G & 79.5 \\
			EfficientNetV2-B3 & IN-22K & 14.4M & 3.0G & 82.5 \\
			EfficientNetV2-S & IN-22K & 21.5M & 8.4G & 84.9 \\
			EfficientNetV2-M & IN-22K & 54.1M & 24.7G & 86.2 \\
			%EfficientNetV2-L & IN-22K & 119.5M & 56.3G & 86.9 \\
			\hline
			\multicolumn{5}{c}{\vspace{2.5em}}\\
	\end{tabular}}
	\caption{Low FLOPs pitfall. EfficientNetV2 results on ImageNet-1K validation set. The red and blue lines denote ImageNet-21K and ImageNet-1K pretraining, respectively.}
	\label{fig:enetv2}
\end{figure}

%%%%%%%%% Method
\section{Approach}
In this section, we investigate the low-FLOPs networks under large-scale pretraining setting. 

\subsection{Architecture: Transformer vs. CNN}
Here we do not propose a new architecture and select the most suitable low-FLOPs network architecture for large-scale visual pretraining.
ViT~\cite{vit} and its variants~\cite{swin,tnt,pvt} have shown the superiority of transformer over CNN in the field of large vision models. As shown in Appendix~\ref{sec:paperswithcode}, Transformer-based models consistently outperform CNNs with similar computational cost when the FLOPs are higher than 5G FLOPs. As for smaller models especially mobile-level model within 600M FLOPs, CNN with inductive bias including locality and s translation equivariance remain dominant. To build efficient backbones for visual tasks, we select CNN as the base model. GhostNet~\cite{ghostnet} is the representative state-of-art mobile model which introduces cheap operation to simplify the standard convolutional layer. 


\subsection{Parameters Are All You Need}
The number of parameters and FLOPs are highly corelated in neural networks. The model with large number of parameters usually has high FLOPs. Considering the intuition that large data requires more parameters, we construct ParameterNet by adding parameters while maintaining low FLOPs.

We start from the conventional convolutional layer. Given the input feature $X\in\mathbb{R}^{C_{in}\times H\times W}$ and the weight tensor $W\in\mathbb{R}^{C_{out}\times C_{in}\times K\times K}$, the conventional convolutional layer operates as
\begin{equation}
	Y = X * W,
\end{equation}
where $Y\in\mathbb{R}^{C_{out}\times H'\times W'}$ is the output, $*$ is the convolution operation and the bias term is omitted for concision. The fully-connected layer can be viewed as the convolutional layer with $1\times1$ kernel size.

Our design principle is adding more parameters while maintaining low FLOPs. There are various approaches to construct ParameterNet, such as dynamic convolution~\cite{dynamic-conv} and re-parameterized convolution~\cite{repvgg}. Although the re-parameterized convolution increase the number of parameters during training, its parameters and FLOPs are unchanged for inference, that is, the model capacity is not increased. In this paper, we mainly consider the efficient dynamic convolution which manyfold increase the number of parameters while almost does not bring in extra FLOPs. 

The dynamic convolution~\cite{dynamic-conv} with $M$ dynamic experts can be written as
\begin{equation}
	\begin{aligned}
		Y &= X * W',\\
		W' &= \sum_{i=1}^{M}\alpha_iW_i.
	\end{aligned}
\end{equation}
where $W_i\in\mathbb{R}^{C_{out}\times C_{in}\times H\times W}$ is the $i$-th convolutional weight tensor and $\alpha_i$ is the corresponding dynamic coefficient. The coefficient $\alpha_i$ is dynamically generated \wrt different input samples, and a typical manner is generating based on the input using MLP module. For the input $X$, a global average pooling is applied to fuse the information into a vector and then a two-layer MLP module with softmax activation is used to produce the coefficients dynamically:
\begin{equation}
	\alpha = \textit{softmax}(\textit{MLP}(\textit{Pool}(X))),
\end{equation}
where $\alpha\in\mathbb{R}^{M}$. 

\paragraph{Complexity Analysis.}
For the standard convolutional layer, the number of parameters is $C_{out}\cdot C_{in}\cdot K\cdot K$ and the number of FLOPs are $H'\cdot W'\cdot C_{out}\cdot C_{in}\cdot K\cdot K$. The dynamic convolution consists of coefficient generation module, dynamic weight fusion and convolution process.
The coefficient generation module with $C_{in}$ hidden dimensions requires $C_{in}^2+C_{in}M$ parameters and $C_{in}^2+C_{in}M$ FLOPs. 
The dynamic weight fusion is parameter-free and has $M\cdot C_{out}\cdot C_{in}\cdot K\cdot K$ FLOPs. 
Thus, the total numbers of parameters and FLOPs of dynamic convolution are $C_{in}^2+C_{in}M + M\cdot C_{out}\cdot C_{in}\cdot K\cdot K$ and $C_{in}^2+C_{in}M + M\cdot C_{out}\cdot C_{in}\cdot K\cdot K+H'\cdot W'\cdot C_{out}\cdot C_{in}\cdot K\cdot K$ respectively. The parameter ratio of dynamic convolution over standard convolution is
\begin{equation}
	\begin{aligned}		
		R_{param} &= \frac{C_{in}^2+C_{in}M + M\cdot C_{out}\cdot C_{in}\cdot K\cdot K}{C_{out}\cdot C_{in}\cdot K\cdot K} \\
		&= \frac{C_{in}}{C_{out}\cdot K\cdot K} + \frac{M}{C_{out}\cdot K\cdot K} + M\\
		&\approx \frac{1}{K^2} + M. \qquad (M<<C_{out}\cdot K\cdot K,~C_{in}\approx C_{out})
	\end{aligned}
\end{equation}
The FLOPs ratio is
\begin{equation}
	\begin{aligned}
		R_{flops} &= \frac{C_{in}^2+C_{in}M + M\cdot C_{out}\cdot C_{in}\cdot K\cdot K+H'\cdot W'\cdot C_{out}\cdot C_{in}\cdot K\cdot K}{H'\cdot W'\cdot C_{out}\cdot C_{in}\cdot K\cdot K} \\
		&= \frac{C_{in}}{H'\cdot W'\cdot C_{out}\cdot K\cdot K} + \frac{M}{H'\cdot W'\cdot C_{out}\cdot K\cdot K} + \frac{M}{H'\cdot W'} + 1\\
		&\approx 1.  \qquad (1<M<<H'\cdot W',~C_{in}\approx C_{out})
	\end{aligned}
\end{equation}
Thus, compared to the standard convolution, the dynamic convolution has about $M\times$ parameters with negligible extra FLOPs. 

%%%%%%%%% Experiment
\section{Experiment}
\subsection{Experimental Settings}
\paragraph{Datasets.}
We adopt the widely-used ImageNet-22K for large-scale pretraing and ImageNet-1K as the normal training data for comparison. ImageNet-22K~\cite{imagenet} is a large-scale image dataset with 14,197,122 images belonging to 21841 categories. ImageNet-1K~\cite{imagenet} is a subset of ImageNet-22K with 1,000 object classes. It contains 1,281,167 training images, and 50,000 validation images.

\begin{table}[htb]
	\vspace{-1.5em}
	\centering
	\small
	\caption{Training hyper-parameters.}
	\label{table-hyper}
	\setlength{\tabcolsep}{8pt}
	\begin{tabular}{l|ccc}
		\hline
		Config & ImageNet-1K & ImageNet-22K & Finetuning \\ 
		\hline
		Epochs & 300 & 90 & 30 \\
		Optimizer & AdamW~\cite{adamw} & AdamW~\cite{adamw} & AdamW~\cite{adamw} \\
		Batch size & 1024 & 4096 & 512 \\
		Start learning rate (LR) & {1e-3} & 4e-3 & 5e-4 \\
		Layer decay & \ding{55} & \ding{55} & 0.5 \\
		Learning rate schedule & {Cosine} & Cosine & Cosine \\
		Warmup epochs & {20} & 5 & 0 \\
		Weight decay & {0.05} & 0.05 & 1e-8 \\
		Label smoothing~\cite{label-smooth} & {0.1} & {0.1} & 0.1 \\
		Stochastic path~\cite{huang2016deep} & \ding{55} & \ding{55} & \ding{55} \\
		RandAugment~\cite{randaugment} & {$\checkmark$} & {$\checkmark$} & $\checkmark$ \\
		Mixup prob.~\cite{mixup} & \ding{55} & \ding{55} & \ding{55} \\
		Cutmix prob.~\cite{cutmix} & \ding{55} & \ding{55} & \ding{55} \\
		Random erasing prob.~\cite{erasing} & {0.25} & {0.25} & \ding{55} \\
		Exponential moving average & {0.9999} & \ding{55} & 0.9999 \\
		\hline
	\end{tabular}
	\vspace{-0.5em}
\end{table}

\paragraph{Training on ImageNet-1K.}
Following the common setting~\cite{deit,swin,convnext}, we train the models for 300 epochs
with 20 warm-up epochs using AdamW~\cite{adamw} optimizer. We use a batch size of 1024. The base learning rate is set as $0.001$ and decays with the cosine schedule. The data augmentation strategy includes RandAugment~\cite{randaugment} and random erasing~\cite{erasing}. Weight decay and label smoothing are adopted for regularization. More details can be found in Table~\ref{table-hyper}.

\paragraph{Pretraining on ImageNet-22K.}
The models are pretraining on ImageNet-22K for 90 epochs with 5 warm-up epochs. The batch size is $4096$ and The base learning rate is set as $0.004$. Other settings basically follow those on ImageNet-1K as shown in Table~\ref{table-hyper}.

\paragraph{Finetuning on ImageNet-1K.}
We finetune the pretrained models on ImageNet-1K for 30 epochs without warm-up epochs. The batch size is $512$ and The base learning rate is set as $0.0005$. The weight decay is set as 1e-8 and random erasing~\cite{erasing} is switched off for better fitting on ImageNet-1K. Other settings basically follow those on ImageNet-1K as shown in Table~\ref{table-hyper}.

\subsection{Main Results}
We build the baseline GhostNet with different FLOPs (\ie, $\sim$300M and $\sim$600M) by tuning the width and depth. Our ParameterNet is constructed by replacing the conventional convolutional layer with dynamic convolution. The number of experts is set as 4 by default. The details of the network architectures are available in the appendix. The results are shown in Table~\ref{tab:ghostnet}. Training only on ImageNet-1K, ParameterNet outperforms the original GhostNet by 0.4-xx accuracy. For GhostNet, pretraining on ImageNet-22K does not help to the performance. ImageNet-22K pretrained ParameterNet can achieve more than 2\% improvement over ImageNet-1K. This indicates that our ParameterNet with more parameters yet similar FLOPs can benefit from the large-scale visual pretraining.

\begin{table*}[htb]
	\centering
	\renewcommand\arraystretch{1.05}
	\caption{ParameterNet results on ImageNet-1K val set.}
	\label{tab:ghostnet}
	\begin{tabular}{l|c|ccc}
		\hline
		& Pretrain data  & Parameters & FLOPs & Top-1  \\
		\hline
		GhostNet-300M~\cite{ghostnet} & IN1K &  8.6M & 289M & 76.3 \\
		GhostNet-300M~\cite{ghostnet} & IN21K &  8.6M & 289M & 75.7 (\textbf{\color{red}-0.6}) \\
		ParameterNet-300M & IN1K &  15.7M & 298M & 77.2 \\
		ParameterNet-300M & IN21K &  15.7M & 298M & 79.1 (\textbf{\color{green}+1.9}) \\
		\hline
		GhostNet-600M~\cite{ghostnet} & IN1K &  19.8M & 579M & 78.6 \\
		GhostNet-600M~\cite{ghostnet} & IN21K &  19.8M & 579M & 78.3 (\textbf{\color{red}-0.3}) \\
		ParameterNet-600M & IN1K &  34.5M & 599M & 79.0 \\
		ParameterNet-600M & IN21K &  34.5M & 599M & 81.6 (\textbf{\color{green}+2.6}) \\
		\hline
	\end{tabular}	
\end{table*}

\paragraph{Comparison with SOTA.}
We compare ParameterNet with other representative models pretrained on ImageNet-22K or larger datasets such as JFT-300M~\cite{sun2017revisiting} and IG-1B-Targeted~\cite{yalniz2019billion}. From the results in Table~\ref{tab:sota}, we can see that our ParameterNet with fewer FLOPs outperforms other models pretrained on large-scale datasets. For example, ParameterNet-600M achieves 81.6\% top-1 accuracy whose \#FLOPs is about 7$\times$ lower than that of ResNet50 or Swin-T.

\begin{table*}[htb]
	\centering
	\renewcommand\arraystretch{1.05}
	\caption{Comparison of ParameterNet and other SOTA models on ImageNet-1K val set.}
	\label{tab:sota}	
	\setlength{\tabcolsep}{4pt}{
	\begin{tabular}{l|c|ccc}
		\hline
		Model & Pretrain data  & Parameters & FLOPs & Top-1  \\
		\hline
		EfficientNet-B0~\cite{noisy-student} & JFT-300M & 5.3M & 390M & 78.1 \\
		Swin-300M~\cite{swin} & IN-22K & 1.9M & 312M & 68.6 \\
		GhostNet-300M~\cite{ghostnet} & IN-22K &  8.6M & 289M & 75.7 \\
		ParameterNet-300M & IN-22K &  15.7M & 298M & \textbf{79.1} \\
		\hline
		ResNet101~\cite{sun2017revisiting} & JFT-300M & 44.5M & 7.8G & 79.2 \\
		ResNet50 (Billion-scale)~\cite{yalniz2019billion} & IG-1B-Targeted & 25.6M & 4.1G & 81.2 \\
		ResNet50 (BiT)~\cite{bit} & IN-22K & 25.6M & 12.0G & 80.0 \\
		EfficientNetV2-B0~\cite{efficientnetv2} & IN-22K &  7.1M & 0.72G & 77.6 \\
		EfficientNetV2-B1~\cite{efficientnetv2} & IN-22K &  8.1M & 1.2G & 79.0 \\
		Swin-T~\cite{swin} & IN-22K & 28M & 4.5G & 80.9 \\
		GhostNet-600M~\cite{ghostnet} & IN-22K &  19.8M & 579M & 78.3 \\
		ParameterNet-600M & IN-22K &  34.5M & 599M & \textbf{81.6} \\
		\hline
	\end{tabular}
	}	
\end{table*}

\paragraph{Inference speed.}
We evaluate the inference speed of ParameterNet and other representative models for comparison. We run models using ONNX toolkit on Intel Xeon Platinum 8378C CPU with single-thread mode. As shown in Figure~\ref{fig:latency}, our ParameterNet outperforms the widely-used ResNet and Swin Transformer for much better accuracy-latency trade-off.

\begin{figure}[tp]
	\centering	
	\setlength{\tabcolsep}{4pt}{
		\begin{tabular}{cc}
			\makecell*[c]{\includegraphics[width=0.5\linewidth]{figures/acc-flops.pdf}}  & \makecell*[c]{\includegraphics[width=0.5\linewidth]{figures/acc-latency.pdf}} 
			\\
			\small (a) Acc v.s. FLOPs & \small (b) Acc v.s. Latency
		\end{tabular}
	}
	\caption{Performance comparison of the representative visual backbone networks with ImageNet-22K pretraining.}
	\label{fig:latency}
\end{figure}


\subsection{Ablation Study}
\paragraph{The number of dynamic experts.}
The number of dynamic experts is an important hyperparameter of dynamic convolution, which directly controls the parameters and FLOPs. As shown in Table~\ref{tab:expert}, more experts will largely increase the number of parameters and slightly influence FLOPs. The performance of more experts improves over fewer experts. We use 4 experts by default for efficiency trade-off.

\begin{table*}[htb]
	\centering
	\renewcommand\arraystretch{1.05}
	\caption{ImageNet-1K val set results \wrt~\#Expert. The base network architecture is GhostNet-300M.}
	\label{tab:expert}
	\begin{tabular}{c|c|ccc}
		\hline
		\#Expert & Pretrain data  & Parameters & FLOPs & Top-1  \\
		\hline
		1 & IN-1K &  8.6M & 289M & 76.3 \\
		1 & IN-22K &  8.6M & 289M & 75.7 \\
		4 & IN-1K &  15.7M & 298M & 77.2 \\
		4 & IN-22K &  15.7M & 298M & 79.1 \\
		8 & IN-1K &  25.2M & 308M & 77.7 \\
		8 & IN-22K &  25.2M & 308M & 79.7 \\
		\hline
	\end{tabular}	
\end{table*}

\paragraph{Dynamic convolution \emph{vs.} re-parameterized convolution.}
As we discussed before, there are various approaches to construct ParameterNet, such as dynamic convolution~\cite{dynamic-conv} and re-parameterized convolution~\cite{repvgg}. We compare these two approaches where the dynamic convolution has 4 experts and the re-parameterized convolution has 3 more paralleled branches based on the original convolution. From the results in Table~\ref{tab:rep}, although the re-parameterized convolution increase the training parameters, its parameters and FLOPs are unchanged for inference, that is, the model capacity is not increased and the ImageNet-22K pretrained performance does not improve.

\begin{table*}[htb]
	\centering
	\renewcommand\arraystretch{1.05}
	\caption{ImageNet-1K val set results \wrt~\#Expert. The base network architecture is GhostNet-300M.}
	\label{tab:rep}
	\begin{tabular}{c|c|cccc}
		\hline
		 & Pretrain  & Train  & Inference & \multirow{2}{*}{FLOPs} & \multirow{2}{*}{Top-1}  \\
		 & data & Parameters & Parameters & & \\
		\hline
		RepConv & IN-1K &  15.7M & 8.6M & 289M & 76.8 \\
		RepConv & IN-22K &  15.7M & 8.6M & 289M & 76.9 \\
		DynamicConv & IN-1K & 15.7M & 15.7M & 298M & 77.2 \\
		DynamicConv & IN-22K & 15.7M & 15.7M & 298M & 79.1 \\
		\hline
	\end{tabular}	
\end{table*}


\paragraph{ParameterNet for other network architectures.}
In addition to CNN, we extend ParameterNet to the transformer architecture (\ie, Swin Transformer). To construct a smaller version, we set the token dimension of Swin-T to 24 to obtain Swin-300M with about 300M FLOPs. From the results in Table~\ref{tab:swin}, the original Swin-300M has a significant accuracy drop when pretraining on ImageNet-22K. Our strategy can achieve +2.2\% performance gain from ImageNet-22K pretraining.

\begin{table*}[htb]
	\centering
	\renewcommand\arraystretch{1.05}
	\caption{ImageNet-1K val set results on Swin-300M architecture.}
	\label{tab:swin}
	\begin{tabular}{l|c|ccc}
		\hline
		 & Pretrain data  & Parameters & FLOPs & Top-1  \\
		\hline
		Original & IN-1K &  1.9M & 312M & 69.2 \\
		Original & IN-22K &  1.9M & 312M & 67.0 (\textbf{\color{red}-2.2}) \\
		Ours & IN-1K &  6.8M & 323M & 72.3 \\
		Ours & IN-22K &  6.8M & 323M & 74.5 (\textbf{\color{green}+2.2}) \\
		\hline
	\end{tabular}	
\end{table*}


\section{Conclusion}
In this paper, we propose a design principle (\ie, ParameterNet) for large-scale visual pretraining by adding more parameters while maintaining low FLOPs.
ParameterNet is a general scheme and has various approaches to implement such as dynamic convolution and re-parameterized convolution. We use the dynamic convolution in practice to construct the ParameterNet models. ParameterNet can overcome the \emph{low FLOPs pitfall} and much benefit from large-scale visual pretraining.
Extensive experiments on ImageNet-22K large-scale dataset have verify the effectiveness of the proposed ParameterNet.
We hope our work can motivate and inspire the future research on vision models, multimodality models and even large language models (LLM).


\begin{appendices}
\sectionfont{\large}
\section{Network Architecture}
The network architecture of GhostNet-300M is the 1.45$\times$ width version~\cite{ghostnet}, and ParameterNet-300M replaces the convolutional layers with dynamic convolutions. The network architecture of ParameterNet-600M is shown in Table~\ref{tab:GhostNet} which is basically built by Ghost blocks. The GhostNet-600M is constructed by replacing dynamic convolutions in ParameterNet-600M with conventional convolutions.

\begin{table}[h]
	\centering
	\renewcommand\arraystretch{1.05}
	\caption{Network architecture of Swin-300M.}
	\label{tab:swin-arch}
	\small
	\begin{tabular}{c|c|c}
		\hline
		Model & $C$ & layer numbers \\ 
		\hline
		Swin-T   & 96  & \{2, 2, 6, 2\} \\
		Swin-300M   & 24  & \{2, 2, 6, 2\} \\
		\hline
	\end{tabular}
\end{table}

\begin{figure}[htp]
	\centering
	% \fbox{\rule{0pt}{2in} \rule{0.9\linewidth}{0pt}}
	\includegraphics[width=0.7\linewidth]{figures/22k_pareto.pdf}	
	\caption{Transformer vs. CNN results on ImageNet-1K validation set with ImageNet-22K pretraining. The data points are from Paperswithcode website in June 2023.}
	\label{fig:pareto}
\end{figure}

\begin{table}[h]
	\centering
	\renewcommand\arraystretch{1.05}
	\caption{Network architecture of ParameterNet-600M. G-bneck denotes Ghost bottleneck with dynamic convolution. The activation function is HSwish~\cite{mobilenetv3}. \#exp means the expansion size. \#out means the number of output channels. SE denotes whether using SE module~\cite{senet}.}
	\label{tab:GhostNet}
	\small
	\begin{tabular}{c|c|c|c|c|c}
		\hline
		Input & Operator & \#exp & \#out & SE & Stride \\ 
		\hline
		$224^2\times 3$   & Conv2d 3$\times$3 &  -  & 32     &  -  & 2 \\
		\hline
		$112^2\times 16$   & G-bneck &  32  & 32     &  -  & 1 \\
		$112^2\times 16$   & G-bneck &  92  & 44     &  -  & 2 \\
		\hline
		$56^2\times 24$   & G-bneck &  136  & 44     &  -  & 1 \\
		$56^2\times 24$   & G-bneck &  136  & 76      &  1  & 2 \\
		\hline
		$28^2\times 40$   & G-bneck &  228  & 76     &  1  & 1 \\
		$28^2\times 40$   & G-bneck &  228  & 76     &  1  & 1 \\
		$28^2\times 40$   & G-bneck &  456  & 152     &  -  & 2 \\
		\hline
		$14^2\times 80$   & G-bneck &  380  & 152     &  -  & 1 \\
		$14^2\times 80$   & G-bneck &  380  & 152    &  -  & 1 \\
		$14^2\times 80$   & G-bneck &  380  & 152    &  -  & 1 \\
		$14^2\times 80$   & G-bneck &  912  & 212     &  1  & 1 \\
		$14^2\times 112$   & G-bneck &  1276  & 212     &  1  & 1 \\
		$14^2\times 112$   & G-bneck &  1276  & 212     &  1  & 1 \\
		$14^2\times 112$   & G-bneck &  1276  & 304     &  1  & 2 \\
		\hline
		$7^2\times 160$   & G-bneck &  1824  & 304      &  1  & 1 \\
		$7^2\times 160$   & G-bneck &  1824  & 304      &  1  & 1 \\
		$7^2\times 160$   & G-bneck &  1824  & 304      &  1  & 1 \\
		$7^2\times 160$   & G-bneck &  1824  & 304      &  1  & 1 \\
		$7^2\times 160$   & G-bneck &  1824  & 304      &  1  & 1 \\
		$7^2\times 160$   & Conv2d 1$\times$1 &  - & 1824     &  -  & 1 \\
		\hline
		$7^2\times 960$   & AvgPool 7$\times$7 &  -  & -      &  -  & - \\\hline
		$1^2\times 960$   & Conv2d 1$\times$1 &  -  & 1280     &  -  & 1 \\\hline
		$1^2\times 1280$   & FC &  -  & 1000     &  -  & - \\
		\hline
	\end{tabular}
\end{table}

For Swin Transformer, we build a smaller version with 300M FLOPs. Compared to the standard Swin-T~\cite{swin}, Swin-300M only changes the dimension to 24. The details are shown in Table~\ref{tab:swin-arch}.

\section{Result Statistics from Paperswithcode}\label{sec:paperswithcode}
We collect the ImageNet-1K finetuning results of ImageNet-22K pretrained models on Paperswithcode website\footnotetext{\url{https://paperswithcode.com/sota/image-classification-on-imagenet}}
 in Table~\ref{fig:pareto}. These results can reflect the trend of CNN and Transformer architectures. The transformer is the better choice when the computational resource is sufficient. The CNNs perform better in low FLOPs setting.
\end{appendices}

{\small
\bibliographystyle{ieee_fullname}
\bibliography{ref}
}


\end{document}
